% -------------------------------------------------------------------------
% WBS ID: RES-GEO-ARR
% Deliverable: Barrier Arrangement, Placement and Spatial Configuration
% Status: In progress
% Evidence status: Corridor-linked placement rules established
% Review status: Not reviewed
% Last updated: 2026-07-19
% Dependencies: RES-DAT-SCN, RES-GEO-TAX, RES-GEO-PAR
% Downstream interfaces: RES-GEO-REP, RES-MOD-IBC, RES-NUM-MSH, Software geometry manifest
% -------------------------------------------------------------------------

\section{RES-GEO-ARR --- Barrier Arrangement, Placement and Spatial Configuration}
\label{sec:res-geo-arr}

\subsection{Purpose and Scope}
\label{sec:res-geo-arr-purpose}

This Deliverable defines how active barrier classes are placed inside
the local corridor coordinate system. Arrangement is constrained by
the evidence-based Kamaishi and Sendai corridor polygons, not by an
arbitrary drawing canvas. The deliverable excludes global coastline
redesign and generative placement optimisation.

\subsection{Research Questions}
\label{sec:res-geo-arr-research-questions}

\begin{enumerate}
  \item How is a barrier positioned relative to the corridor
    centreline $x'$ and width direction $y'$?
  \item How are continuous walls distinguished from finite obstacles
    and arrays?
  \item Which placement choices are physical scenario decisions, and
    which are numerical buffer or mesh-quality decisions?
\end{enumerate}

\subsection{Terminology and Definitions}
\label{sec:res-geo-arr-definitions}

$x'$ is the local propagation-direction coordinate supplied by
RES-DAT-SCN, $y'$ is the transverse corridor coordinate, and $z'$ is
vertical elevation in the solver datum. The angle $\psi_j$ is the
plan-view orientation of barrier $j$ measured relative to the
$x'$ axis; $\psi_j=\pi/2$ denotes a barrier aligned across the
corridor.

\subsection{Literature Evidence}
\label{sec:res-geo-arr-literature-evidence}

\textcite{mori2012nationwide} supports retaining different placement
contexts for Sanriku/Kamaishi and Sendai because the same event
produced contrasting run-up and inundation behaviour. Source role:
supports corridor-specific placement rather than one generic test
site. \textcite{watanabe2022validation} supports preserving exact
probe and structure locations for validation-facing local simulations.
Limitation: neither source supplies final barrier coordinates; those
remain RES-DAT-SCN/RES-DAT-CAS data products.

\subsection{Governing Relations and Quantitative Definitions}
\label{sec:res-geo-arr-governing-relations}

The plan-view local placement vector is:

\begin{align}
  \mathbf{r}'_j
  &=
  \left[
    x'_j,\,
    y'_j
  \right]^{\mathsf{T}}
  \label{eq:res-geo-arr-placement-vector}
\end{align}

The local basis for each barrier is:

\begin{align}
  \mathbf{e}_{n,j}
  &=
  \left[
    \cos\psi_j,\,
    \sin\psi_j
  \right]^{\mathsf{T}}
  \label{eq:res-geo-arr-barrier-normal}
  \\
  \mathbf{e}_{t,j}
  &=
  \left[
    -\sin\psi_j,\,
    \cos\psi_j
  \right]^{\mathsf{T}}
  \label{eq:res-geo-arr-barrier-tangent}
\end{align}

Decision role: these definitions make barrier placement reproducible
from the corridor transform. Supporting evidence: RES-DAT-SCN owns
the corridor construction; \textcite{watanabe2022validation} supports
validation-facing control of structure and probe locations.

\subsection{Geometry Family or Representation}
\label{sec:res-geo-arr-geometry-family-representation}

Continuous walls are represented as elongated patches whose tangent
direction is $\mathbf{e}_{t,j}$. Obstacle and array cases are
represented as one or more finite elements with recorded centre
points, spacing and row index.

\subsection{Design Variables and Parameter Bounds}
\label{sec:res-geo-arr-design-variables-parameter-bounds}

Arrangement variables include $x'_j$, $y'_j$, $\psi_j$, number of rows,
along-row spacing, cross-row spacing, stagger offset and gap/opening
fraction. These variables are comparison controls, not yet an
optimisation search space.

\subsection{Geometric Descriptors}
\label{sec:res-geo-arr-geometric-descriptors}

Placement descriptors shall include clearance to artificial inlet,
outlet and side boundaries; clearance to steep terrain changes; and
distance from run-up, overtopping and wake measurement sections.

\subsection{Placement and Arrangement Rules}
\label{sec:res-geo-arr-placement-arrangement-rules}

Baseline placement rules are:

\begin{enumerate}
  \item place wall-type cases approximately transverse to $x'$ unless
    a real coastline, harbour or defence alignment provides a citable
    alternative;
  \item place array-type cases with all elements inside the accepted
    terrain/corridor polygon and away from sponge zones;
  \item keep comparison sections fixed when comparing barrier classes;
  \item do not narrow the corridor to fit a barrier unless the
    narrowing represents a physical bay or coastline constraint.
\end{enumerate}

\subsection{Feasibility Constraints}
\label{sec:res-geo-arr-feasibility-constraints}

A placement is infeasible if it intersects missing terrain data,
requires a nonphysical wall at the side boundary, violates local mesh
quality constraints or places the barrier inside a numerical
relaxation region.

\subsection{Two- and Three-Dimensional Representation}
\label{sec:res-geo-arr-two-three-dimensional-representation}

% TODO: Complete this RES-GEO Domain-specific subsection.

\subsection{Candidate Approaches}
\label{sec:res-geo-arr-candidate-approaches}

% TODO: Define the credible candidate approaches considered.

\subsection{Critical Comparison}
\label{sec:res-geo-arr-critical-comparison}

% TODO: Compare candidates using physically and project-relevant criteria.
% Do not select an approach solely on popularity or implementation ease.

\subsection{Assumptions and Applicability}
\label{sec:res-geo-arr-assumptions}

% TODO: State assumptions and the conditions under which they are valid.

\subsection{Limitations and Failure Conditions}
\label{sec:res-geo-arr-limitations}

% TODO: State where the evidence, model or selected approach ceases to be
% reliable or applicable.

\subsection{Project-Specific Conclusions}
\label{sec:res-geo-arr-conclusions}

For the Studentship campaign, arrangement is part of the data/model
interface. It links source-to-coast trajectory, terrain extraction,
barrier geometry and validation probes; it is not a separate design
art exercise.

\subsection{Selected Approach and Rationale}
\label{sec:res-geo-arr-selected-approach}

The selected approach is corridor-relative placement using
$(x',y',z')$ coordinates and fixed comparison sections. Free placement
optimisation is deferred.

\subsection{Unresolved Decisions}
\label{sec:res-geo-arr-unresolved-decisions}

Open decisions are final Kamaishi and Sendai target coordinates, final
comparison sections, sponge widths, and any real-defence alignment
source for the Kamaishi breakwater.

\subsection{Requirements and Handoffs}
\label{sec:res-geo-arr-requirements}

Software shall generate barrier placement from corridor coordinates,
clip it against accepted terrain/corridor polygons, and emit placement
diagnostics for mesh quality and validation probe locations.

\subsection{Deliverable Completion Record}
\label{sec:res-geo-arr-completion-record}

% TODO: Record whether the Deliverable is:
%
% - Not started
% - In progress
% - Draft complete
% - Reviewed
% - Accepted
%
% Acceptance requires:
%
% - the research questions to be answered;
% - claims to be supported by appropriate evidence;
% - assumptions and limitations to be explicit;
% - the project-specific conclusion to be stated;
% - downstream requirements to be traceable;
% - unresolved decisions to be closed or formally deferred.
